\documentclass[10pt]{article}
\usepackage[margin=1in]{geometry}
\usepackage{amsmath}
\usepackage{booktabs}
\usepackage{graphicx}
\usepackage{hyperref}

\title{Linguistic Entropy Coding: Reversible Ciphertext Embedding Through
Language-Model Sampling}
\author{Anonymous authors}
\date{}

\begin{document}
\maketitle

\begin{abstract}
We study a reversible mapping in which authenticated ciphertext substitutes
for the random entropy ordinarily consumed during language-model sampling.
An arithmetic interval coder selects tokens from deterministically quantized
next-token distributions, allowing an unchanged token sequence to recover the
exact ciphertext. We present a minimal reproducible implementation and evaluate
recovery, capacity, linguistic quality, and detectability relative to ordinary
sampling. This work does not propose a new cryptographic primitive or claim
novelty for neural linguistic steganography.
\end{abstract}

\section{Introduction}
For an autoregressive model,
\[
  P(x_{1:T})=\prod_{t=1}^{T}P(x_t\mid x_{<t}),
\]
ordinary stochastic generation consumes random entropy to select each token.
LEC replaces that entropy source with authenticated ciphertext, which should be
computationally indistinguishable from random bits under standard assumptions.

\section{Related Work}
\paragraph{Linguistic steganography.}
Discuss classical and neural linguistic steganography, distribution-preserving
coding, and threat models.
\paragraph{Entropy coding.}
Discuss arithmetic coding, range coding, and asymmetric numeral systems.
\paragraph{Language-model sampling.}
Discuss temperature and candidate truncation.

\section{Linguistic Entropy Coding}
At step $t$, floating-point model probabilities are converted to integer
frequencies $F_{t,i}$ satisfying
\[
  \sum_i F_{t,i}=2^N.
\]
The encrypted payload identifies an integer point in an initial interval. Each
observed token narrows the interval according to $F_t$ until it contains one
point. The decoder replays the same distributions and interval updates.

\section{Implementation}
The v0.1 implementation fixes DistilGPT-2 weights and tokenizer, deterministic
CPU inference, one prompt, ChaCha20-Poly1305, fixed payload blocks, and
tokenization-roundtrip-safe candidates. Integer largest-remainder quantization
uses token identifiers as deterministic tie breakers.

\section{Experiments}
\subsection{RQ1: Reversibility}
Report byte recovery accuracy over payload sizes and configurations.
\subsection{RQ2: Capacity}
Report bits per token and entropy utilization
\[
  \eta=\frac{\text{embedded bits}}{\sum_t H(P_t)}.
\]
\subsection{RQ3--RQ4: Quality and Tradeoff}
Compare ordinary and LEC sampling with independent-model perplexity and blinded
preference judgments. Plot capacity against quality for temperature, top-$k$,
and quantization settings.
\subsection{RQ5: Detectability}
Train and evaluate classifiers on held-out ordinary versus LEC completions.
Report ROC--AUC with confidence intervals and corpus-level splits.

\section{Results}
\begin{figure}[h]
  \centering
  \fbox{\parbox[c][1.6in][c]{0.75\linewidth}{\centering
  Capacity--quality Pareto curve placeholder}}
  \caption{Embedding capacity versus linguistic quality.}
\end{figure}

\section{Limitations}
Decoding requires an exact model, tokenizer, runtime configuration, and
unmodified text. The method supplies no paraphrase resistance, error correction,
or production key management. Model-generated text may remain detectable.

\section{Discussion}
Future work should test whether ciphertext can serve as reusable sampling
entropy in other probabilistic generative models and examine finite-precision
and synchronization costs.

\section{Conclusion}
Information normally consumed as random sampling entropy can carry encrypted
data while following quantized language-model distributions.

\end{document}
